% Source-preserving cumulative Indonesian driver for the substantial Chapter 13 boundary.
% The staging directory supplies frozen Volume 1 dependencies, preserves the admitted
% S111-S132 target bytes, and overlays mt13 plus the complete S133-S136 target sources.
\input volwp.aux
\mtlulustyle
\luluvolumeno=1

% Localize reader-facing headings generated by frozen source macros.
\def\Notesheader#1{\immediate\write0{#1 Catatan dan komentar}
     \wheader{#1 Catatan}{16}{8}{5}{48pt}
     \noindent{\bf #1 Catatan dan komentar}}
\def\vNotesheader#1#2{\immediate\write0{#2 Catatan dan komentar}
     \wheader{#2 Catatan}{16}{8}{5}{#1}
     \noindent{\bf #2 Catatan dan komentar}}
\def\Hint#1{({\it Petunjuk\/}: #1)}
\long\def\proof#1{\ifwithproofs{\medskip

     \noindent{\bf Bukti} #1}\else\fi}

% Preserve the four source diagrams used by Section 113.
\input mt113-dvipdfmx-images
% Preserve the three source diagrams used by Section 134.
\input mt134-dvipdfmx-images

\special{pdf:docinfo <<
  /Title (Fondasi Teori Ukuran - Volume 1, Bab 13 lengkap sampai Bagian 136)
  /Author (D. H. Fremlin; adaptasi Bahasa Indonesia atas arahan pengguna)
  /Subject (Adaptasi Bahasa Indonesia dari Measure Theory, Volume 1, halaman resmi 10-90)
  /Keywords (teori ukuran, integral, ukuran Lebesgue, integral Riemann, himpunan Cantor, garis real diperluas, teorema kelas monoton, id-ID, O007)
>>}
\begingroup
\catcode64=12
\special{pdf:put @catalog << /Lang (id-ID) /ViewerPreferences << /DisplayDocTitle true >> >>}
\endgroup

\pageno=1
\vbox to \vsize{
  \vskip 42pt
  \centerline{\fourteenbf FONDASI TEORI UKURAN}
  \vskip 15pt
  \centerline{\bf Volume 1: Minimum yang Tak Tereduksi}
  \vskip 34pt
  \centerline{\fourteenbf Batas kumulatif Bab 13}
  \vskip 8pt
  \centerline{\fourteenbf Halaman resmi 10--90}
  \vskip 38pt
  \centerline{D. H. Fremlin}
  \vskip 6pt
  \centerline{Adaptasi Bahasa Indonesia atas arahan pengguna}
  \vfill
  \hsize=350pt
  \leftskip=35pt
  \rightskip=35pt
  \noindent Sumber: \it Measure Theory, Volume 1: The Irreducible Minimum,
  \rm Bagian 111--115, 121--123, pendahuluan Bab 13, dan Bagian 131--136,
  D. H. Fremlin. Cakupan unik sumber resmi: halaman 10--90 (81 halaman).
  Terjemahan dan modernisasi pembaca: 23 Agustus 2026. Materi turunan
  Fremlin tetap berada di bawah Design Science License. Sumber yang dapat
  diedit, lisensi, atribusi, backend semantik, dan bukti QA disertakan.
  \vskip 10pt
  \noindent Provenans produksi berbantuan model: OpenAI Codex
  gpt-5.6-sol, Ultra. Seluruh kredit penulis dan kontributor dipertahankan.
  \vskip 18pt
}
\eject

\pageno=1
\def\folio{\number\pageno}
\def\readerheadline#1#2{\headline={\rlap{\eightbf #1}\hfill{\eightit #2}\hfill
  \llap{\eightrm\folio}}}
\readerheadline{111}{Aljabar-$\sigma$}
\footline={\hfill\eightrm\folio\hfill}
\input mt111
\eject
\readerheadline{112}{Ruang ukur}
\input mt112
\eject
\readerheadline{113}{Ukuran luar dan konstruksi Carath\'eodory}
\input mt113
\eject
\readerheadline{114}{Ukuran Lebesgue pada $\Bbb R$}
\input mt114
\eject
\readerheadline{115}{Ukuran Lebesgue pada $\Bbb R^r$}
\input mt115
\eject
\readerheadline{121}{Fungsi terukur}
\input mt121
\eject
\readerheadline{122}{Definisi integral}
\input mt122
\eject
\readerheadline{123}{Teorema-teorema konvergensi}
\input mt123
\eject

% Source order: the Chapter 13 introduction is official printed page 57.
\readerheadline{Bab 13}{Pelengkap}
% The frozen Fremlin macro prints an English chapter label and overwrites the
% localized introduction heading.  Override presentation only for this id-ID
% reader while preserving the source chapter number, name, and topology.
\def\newchapter#1{\gdef\topparagraph{}
   \gdef\bottomparagraph{Bab #1 {\it pendahuluan}}
   \gdef\newparagraph{Bab #1 {\it pendahuluan}}
   \footmarkcount=0
   \gdef\sectionname{Pendahuluan}
   \gdef\headlinesectionname{Pendahuluan}
   \centerline{\bf Bab #1}
   \medskip
   \centerline{\bf \chaptername}
   \medskip
   }
\input mt13
\eject
\readerheadline{131}{Subruang terukur}
\input mt131
\eject
\readerheadline{132}{Ukuran luar dari ukuran}
\input mt132
\eject
\readerheadline{133}{Konsep integrasi yang lebih luas}
\input mt133
\eject
\readerheadline{134}{Lebih lanjut tentang ukuran Lebesgue}
\input mt134
\eject
\readerheadline{135}{Garis real diperluas}
\input mt135
\eject
\readerheadline{136}{Teorema Kelas Monoton}
\input mt136
\end
